% Figure 1: Six-axis design-space diagram for §3.
% Six axes radiate from a central origin, each labeled with its name
% and a representative range of values. Selected configurations from
% §5 are placed as labeled markers along the axes their values pin.
\begin{figure}[t]
  \centering
  \begin{tikzpicture}[
      axis/.style={->, thick, gray!70!black},
      axislabel/.style={font=\footnotesize\bfseries, align=center},
      valuelabel/.style={font=\scriptsize, gray!50!black, align=center},
      cfg/.style={
        circle, draw, fill=white, inner sep=1pt,
        font=\scriptsize\bfseries, minimum size=4mm
      },
      depedge/.style={->, dashed, gray!60, thick},
      scale=1.0, every node/.style={transform shape}
    ]
    % Six axes at 30°, 90°, 150°, 210°, 270°, 330° (rotated 30° so
    % no axis lies horizontally — readability)
    \def\axislen{2.6}
    \foreach \angle/\name/\range in {%
       30/{Anchor chain}/{Stellar, Eth L1, L2, Mina, Aztec},%
       90/{SNARK system}/{Groth16, PLONK, Halo2, STARK, Plonky2/3},%
      150/{Curve {\&} hash}/{BLS12-381, BN254, Pasta, Goldilocks},%
      210/{PQ stance}/{Pairing, EC, Hash},%
      270/{Trusted setup}/{Per-circuit, Universal, Transparent},%
      330/{Verifier host}/{Native, Bytecode, Verifier-as-rollup}}{%
        \draw[axis] (0,0) -- (\angle:\axislen);
        \node[axislabel] at (\angle:\axislen+0.65) {\name};
        \node[valuelabel] at (\angle:\axislen+1.35) {\range};
    }
    % Configurations placed near the dimension they most clearly pin.
    % Positions chosen for legibility, not as quantitative claims.
    \node[cfg]               (C1) at ( 30:1.2) {C1};
    \node[cfg]               (C2) at ( 90:1.2) {C2};
    \node[cfg]               (C3) at (150:1.2) {C3};
    \node[cfg]               (C4) at (210:1.2) {C4};
    \node[cfg]               (C5) at (270:1.2) {C5};
    \node[cfg]               (C6) at (330:1.2) {C6};
    \node[cfg, dashed]       (C7) at (  0:0.4) {C7};

    % Origin marker
    \fill[gray!60] (0,0) circle (1.2pt);

    % Dependency arrows from §3.7: chain → host → curve → SNARK → setup.
    % Drawn as a small inset chain at the figure's bottom.
    \begin{scope}[shift={(-3.2,-3.6)}, font=\scriptsize]
      \node[draw, rounded corners=2pt, inner sep=2pt] (C) {chain};
      \node[draw, rounded corners=2pt, right=4mm of C, inner sep=2pt] (H) {host};
      \node[draw, rounded corners=2pt, right=4mm of H, inner sep=2pt] (Cu) {curve};
      \node[draw, rounded corners=2pt, right=4mm of Cu, inner sep=2pt] (S) {SNARK};
      \node[draw, rounded corners=2pt, right=4mm of S, inner sep=2pt] (Se) {setup};
      \draw[depedge] (C) -- (H);
      \draw[depedge] (H) -- (Cu);
      \draw[depedge] (Cu) -- (S);
      \draw[depedge] (S) -- (Se);
      \node[font=\tiny, gray!60!black, above=0.5mm of H, xshift=10mm]
        {dependency chain (\S\ref{sec:dimensions-deps})};
    \end{scope}
  \end{tikzpicture}
  \caption{Six design axes for metadata-hiding group registries with
    SNARK-gated state changes. Configurations C1--C7
    (\S\ref{sec:configurations}) are placed as markers near the axes
    they most clearly pin; placements are illustrative rather than
    quantitative. C7 (dashed) is the host-function-blocked
    hypothetical. The inset chain at the bottom shows the §3.7
    dependency direction: fixing the anchor chain transitively
    constrains four of the six axes.}
  \label{fig:axes}
\end{figure}
